% ─────────────────────────────────────────────────────────────────────────────
% _preamble.tex — Préambule commun à tous les fichiers maîtres du projet.
%
% Usage dans chaque fichier maître :
%   \documentclass[12pt,twoside,openright]{memoir}
%   % optionnel : \def\AVECNOTES{} et/ou \def\AVECILLUSTRATIONS{} et/ou \def\AVECMUSIQUE{}
%   % ─────────────────────────────────────────────────────────────────────────────
% _preamble.tex — Préambule commun à tous les fichiers maîtres du projet.
%
% Usage dans chaque fichier maître :
%   \documentclass[12pt,twoside,openright]{memoir}
%   % optionnel : \def\AVECNOTES{} et/ou \def\AVECILLUSTRATIONS{} et/ou \def\AVECMUSIQUE{}
%   % ─────────────────────────────────────────────────────────────────────────────
% _preamble.tex — Préambule commun à tous les fichiers maîtres du projet.
%
% Usage dans chaque fichier maître :
%   \documentclass[12pt,twoside,openright]{memoir}
%   % optionnel : \def\AVECNOTES{} et/ou \def\AVECILLUSTRATIONS{} et/ou \def\AVECMUSIQUE{}
%   % ─────────────────────────────────────────────────────────────────────────────
% _preamble.tex — Préambule commun à tous les fichiers maîtres du projet.
%
% Usage dans chaque fichier maître :
%   \documentclass[12pt,twoside,openright]{memoir}
%   % optionnel : \def\AVECNOTES{} et/ou \def\AVECILLUSTRATIONS{} et/ou \def\AVECMUSIQUE{}
%   \input{_preamble}
%   \begin{document}
%   ...
% ─────────────────────────────────────────────────────────────────────────────

% ─── PDF VERSION ──────────────────────────────────────────────────────────────
\ifdefined\pdfvariable\pdfvariable minorversion 7\fi

% ─── PACKAGES ─────────────────────────────────────────────────────────────────
\usepackage{fontspec}
% ────── EB Garamond
%%% \setmainfont{EB Garamond}
% ────── Cormorant Garamond
%%% \setmainfont[
%%%   UprightFont    = *-Regular,
%%%   ItalicFont     = *-Italic,
%%%   BoldFont       = *-SemiBold,
%%%   BoldItalicFont = *-SemiBoldItalic
%%% ]{Cormorant Garamond}
% ────── EB Garamond (main) + Adobe Garamond Pro (bold only)
\setmainfont{EB Garamond}[
  BoldFont       = AGaramondPro-Semibold,
  BoldItalicFont = AGaramondPro-SemiboldItalic,
]

\newfontfamily{\arabicfont}{Geeza Pro}
\newcommand{\arab}[1]{{\arabicfont\textdir TRT #1}}
\newfontfamily{\cjkfont}{Hiragino Mincho ProN}
\newcommand{\cjk}[1]{{\cjkfont #1}}
\newcommand{\shouting}[1]{\textbf{\textsc{#1}}}
\usepackage[russian,french]{babel}
\usepackage[program=/usr/local/bin/lilypond,includepaths={music}]{lyluatex}
\usepackage{microtype}
\usepackage{fmtcount}
\usepackage{catchfile}
\usepackage{graphicx}
\usepackage{xcolor}
\usepackage{amsmath}
\usepackage{xparse}
\usepackage{tikz}
\usepackage{afterpage}
\usepackage{needspace}
\usepackage{pdfpages}
\usepackage[absolute,overlay]{textpos}
\usepackage{tabularx}
\usepackage{booktabs}
\usepackage{calc}
\usepackage{wrapfig}
\usepackage{tcolorbox}
\usepackage{ragged2e}
% Fix: \WF@finale restores \linewidth and \@totalleftmargin only when
% \parshape==\WF@fudgeparshape; the "forced to float" collision path bypasses
% that guard, leaving both corrupted globally.  Patch \WF@finale to force-reset
% both after every call, regardless of which path was taken.
\makeatletter
\let\WF@@orig@finale\WF@finale
\def\WF@finale{\WF@@orig@finale\global\linewidth\textwidth\global\@totalleftmargin\z@}
\makeatother
% Safety net: clear any residual wrapfig state at chapter boundaries.
\pretocmd{\chapter}{\WFclear}{}{}
\usepackage{hyperref}
\usepackage{bookmark}
\hypersetup{pdfpagelayout=TwoPageRight, hidelinks}
\usetikzlibrary{calc,positioning}


% ─── VERSION / GIT ────────────────────────────────────────────────────────────
\input{version.tex}
\IfFileExists{tete_de_veau_ravigote.gitinfo}%
  {\CatchFileDef{\gitcommit}{tete_de_veau_ravigote.gitinfo}{\endlinechar=-1}}%
  {\def\gitcommit{}}

% ─── GRAYSCALE IMAGE FILTER ───────────────────────────────────────────────────
\graphicspath{{iconography/}}
\newcommand{\figcaptionname}{\footnotesize}
\newcommand{\figcaptiondetail}{\scriptsize}
\newcommand{\figpagetitlename}{\small\scshape}
\newcommand{\figpagetitledetail}{\footnotesize}
\newcommand{\bwimage}[2][]{%
  \begin{tikzpicture}
    \node[inner sep=0, outer sep=0pt, fill=white] (img) {\includegraphics[#1]{#2}};
    %%% begin{scope}[blend mode=saturation]
    %%%  \fill[black] (img.south west) rectangle (img.north east);
    %%% end{scope}
  \end{tikzpicture}%
}

% ─── ESPACEMENT FLOTTANTS ─────────────────────────────────────────────────────
\setlength{\textfloatsep}{8pt}   % espace entre flottant top/bottom et le texte (défaut ~20pt)

% ─── ICONOGRAPHIE ─────────────────────────────────────────────────────────────
% All commands below are no-ops unless \AVECILLUSTRATIONS is defined.
% Image paths are relative to the graphicspath (iconography/).
\ifdefined\AVECILLUSTRATIONS
  \newlength{\iconofillheight}% scratch register for [p] page-fill placement
  % \iconographieimg[scale]{image}{caption}
  % Single image centred on its own dedicated page (via \afterpage).
  % [scale] height as fraction of \textheight (default 0.75); {caption} may be empty.
  \newcommand{\iconographieimg}[3][0.75]{%
    \afterpage{%
      \thispagestyle{gallimard}%
      \vspace*{\fill}%
      \begin{center}%
        \bwimage[height=#1\textheight,width=\textwidth,keepaspectratio]{#2}\\[0.6em]%
        {\figcaptionname\itshape #3}%
      \end{center}%
      \vspace*{\fill}%
      \newpage
    }%
  }
  % \iconographietex{figures/foo_fig.tex}
  % Inserts one hand-crafted figure page from an external .tex file (via \afterpage).
  % The file is responsible for its own layout — see figures/*_fig.tex for examples.
  \newcommand{\iconographietex}[1]{%
    \afterpage{%
      \thispagestyle{gallimard}%
      \input{#1}%
      \clearpage
    }%
  }
  % \iconographietexpair{figures/foo_fig.tex}{figures/bar_fig.tex}
  % Like \iconographietex but inserts two consecutive figure pages.
  \newcommand{\iconographietexpair}[2]{%
    \afterpage{%
      \thispagestyle{gallimard}%
      \input{#1}%
      \input{#2}%
      \clearpage
    }%
  }
  % \iconographiewrapfig{pos}{width}{\bwimage[width=\linewidth]{path}}{caption}
  % Image wrapped by surrounding text (wrapfigure).
  % {pos} = l or r;  {width} = e.g. 0.35\textwidth;  {caption} may be empty {}.
  \newcommand{\iconographiewrapfig}[4]{%
    \setlength{\intextsep}{0pt}%
    \setlength{\columnsep}{8pt}%
    \begin{wrapfigure}{#1}{#2}%
      \centering
      #3%
      \ifx&#4&\else\\[0.2em]{\figcaptionname\itshape #4}\\[0.4\baselineskip]\fi%
    \end{wrapfigure}%
  }
  % \iconographieinlineblock[placement][lines]{\bwimage[width=W]{path}}{caption}
  % Centred image block in the text flow. {caption} may be empty {}.
  % [placement] empty or omitted → inline: caller must put \par before and \noindent after.
  % [placement] = t or b → LaTeX float (top/bottom of page); no \par/\noindent needed.
  % [placement] = p → fill remaining page height inline (image scaled to fit, no float).
  % [lines] number of caption lines (default 1); used by [p] to reserve the right height.
  % Multiple inline calls in the same paragraph are supported.
  \NewDocumentCommand{\iconographieinlineblock}{O{} O{1} m m}{%
    % #1=placement  #2=caption lines (for [p])  #3=image  #4=caption
    \ifx&#1&%
      % no placement → inline
      \vspace{0.5\onelineskip}%
      \centerline{#3}%
      \ifx&#4&\else\par\centerline{\figcaptionname\itshape#4}\vspace{0.5\baselineskip}\fi%
      \vspace{-0.5\onelineskip}%
    \else
      \ifthenelse{\equal{#1}{p}}{%
        % p → scale image to fill remaining page height
        \par
        \begingroup
          \setlength{\iconofillheight}{\dimexpr\pagegoal-\pagetotal\relax}%
          \addtolength{\iconofillheight}{-2\baselineskip}%
          \ifx&#4&\else\addtolength{\iconofillheight}{-#2\baselineskip-0.5\baselineskip}\fi%
          \ifdim\iconofillheight>0pt
            \vspace{2\baselineskip}%
            \centerline{\resizebox*{!}{\iconofillheight}{#3}}%
            \ifx&#4&\else\par{\centering\figcaptionname\itshape#4\par}\fi%
          \fi
        \endgroup
      }{%
        % placement given → float
        \begin{figure}[#1]
          \centering#3%
          \ifx&#4&\else\par{\figcaptionname\itshape#4}\fi%
        \end{figure}%
      }%
    \fi
  }
  % \iconographiedouble[placement][height][margin]{\bwimage{img1}}{\bwimage{img2}}[caption]
  % Two images side by side at equal height.
  % [placement] = b (default, push to bottom) or t (push to top)
  % [height] common image height (default 5cm); [margin] symmetric indent (default 0pt)
  % [caption] shared caption, omit for none.
  \NewDocumentCommand{\iconographiedouble}{O{b} O{5cm} O{0pt} m m O{}}{%
    \ifthenelse{\equal{#1}{b}}{\vspace*{\fill}}{\vspace*{0pt}}%
    \vspace{0.5\onelineskip}%
    {\setlength{\leftskip}{#3}\setlength{\rightskip}{#3}%
    \noindent\bwimage[height=#2]{#4}\hfill\bwimage[height=#2]{#5}\par}%
    \ifx&#6&%
      \vspace{0.5\onelineskip}%
    \else%
      \vspace{0.2\onelineskip}%
      {\centering\figcaptionname\itshape#6\par}%
      \vspace{0.3\onelineskip}%
    \fi%
    \ifthenelse{\equal{#1}{t}}{\vspace*{\fill}}{}%
  }
  % \iconographietriple[placement][height][margin]{\bwimage{img1}}{\bwimage{img2}}{\bwimage{img3}}[caption]
  % Three images side by side at equal height. Same options as \iconographiedouble.
  \NewDocumentCommand{\iconographietriple}{O{b} O{5cm} O{0pt} m m m O{}}{%
    \ifthenelse{\equal{#1}{b}}{\vspace*{\fill}}{\vspace*{0pt}}%
    \vspace{0.5\onelineskip}%
    {\setlength{\leftskip}{#3}\setlength{\rightskip}{#3}%
    \noindent\bwimage[height=#2]{#4}\hfill\bwimage[height=#2]{#5}\hfill\bwimage[height=#2]{#6}\par}%
    \ifx&#7&%
      \vspace{0.5\onelineskip}%
    \else%
      \vspace{0.2\onelineskip}%
      {\centering\figcaptionname\itshape#7\par}%
      \vspace{0.3\onelineskip}%
    \fi%
    \ifthenelse{\equal{#1}{t}}{\vspace*{\fill}}{}%
  }
\else
  \newcommand{\iconographieimg}[3][0.75]{}
  \newcommand{\iconographietex}[1]{}
  \newcommand{\iconographietexpair}[2]{}
  \newcommand{\iconographiewrapfig}[4]{}
  \NewDocumentCommand{\iconographieinlineblock}{O{} O{1} m m}{}
  \NewDocumentCommand{\iconographiedouble}{O{b} O{5cm} O{0pt} m m O{}}{}
  \NewDocumentCommand{\iconographietriple}{O{b} O{5cm} O{0pt} m m m O{}}{}
\fi

% ─── EXTRAITS MUSICAUX ────────────────────────────────────────────────────────
% Compiler avec : lualatex "\def\AVECMUSIQUE{}\input{...}"
\ifdefined\AVECMUSIQUE
  % Extrait musical inline : [#1]=espacement (défaut 0.66), #2=fichier dans music/, #3=légende
  \newcommand{\score}[3][0.66]{%
    \addcontentsline{mus}{score}{#3}%
    \noindent\begin{minipage}{\textwidth}%
    \vspace{#1\onelineskip}%
    \lilypondfile[line-width=\textwidth]{music/#2}%
    \begin{center}%
      \vspace{-#1\onelineskip}%
      {\centering\figcaptionname\itshape#3}%
      \vspace{#1\onelineskip}%
    \end{center}%
    \end{minipage}}
\else
  \newcommand{\score}[3][1]{}
\fi

% ─── NOTES FINALES ────────────────────────────────────────────────────────────
% Compiler avec : lualatex "\def\AVECNOTES{}\input{...}"
% \ENDOFBOOKNOTES : notes regroupées en fin de livre (avec en-tête de scène) ;
%                   absent = notes à la fin de chaque scène (sans en-tête).
\usepackage{endnotes}
\ifdefined\AVECNOTES
  \newcommand{\nf}[1]{\endnote{#1}}
\else
  \newcommand{\nf}[1]{}
\fi
% Supprimer le titre automatique du package (on gère le nôtre)
\renewcommand{\notesname}{}
% Style des notes : même fonte, corps légèrement réduit
\renewcommand{\enotesize}{\small}
% Format des notes en fin de livre : sans retrait, numéro en corps normal
\makeatletter
\renewcommand\enoteformat{%
  \rightskip\z@ \leftskip\z@ \parindent=0pt
  \leavevmode\theenmark.\enspace}
\makeatother
% Fix PDF ActualText for œ (oe-ligature) copy-paste in \source URLs
\usepackage{accsupp}
\newcommand{\oechar}{\char"153\relax}
{\catcode"153=\active\gdefœ{\BeginAccSupp{method=hex,unicode,ActualText=0153}\oechar\EndAccSupp{}}}
% Source dans les notes : saut de ligne avant le lien Wikipedia
\newcommand{\source}[1]{\newline\textit{Source~:}~{\upshape\addfontfeatures{RawFeature=-tlig;-calt}\def\_{\textunderscore\allowbreak}\def\%{\char"25\relax\allowbreak}\catcode"153=\active\scantokens{#1\empty}}}
% Image seule dans une note (centrée sous le texte)
\newcommand{\nfimg}[2]{\newline\includegraphics[width=#1\linewidth]{#2}}
% Image côte à côte avec le texte : #1=texte, #2=frac image, #3=path, #4=url source
\newcommand{\nfimgblock}[4]{%
  \begin{minipage}[t]{0.45\linewidth}\small#1\end{minipage}%
  \hfill%
  \begin{minipage}[t]{#2\linewidth}\includegraphics[width=\linewidth]{#3}\end{minipage}%
  \source{#4}}
% Filet ornemental (séparateur de scène intra-chapitre)
\newcommand{\ornamentalrule}{%
  \vspace{\onelineskip}%
  {\centering\rule{0.4\textwidth}{0.4pt}\par}%
  \vspace{\onelineskip}%
}
% Font for scene header in notes block
\makeatletter
\ifdefined\ENDOFBOOKNOTES
  \newcommand{\scene@noteheader}[1]{{\normalfont\large #1}}
\else
  \newcommand{\scene@noteheader}[1]{{\normalfont\scshape #1}}
\fi
% \scene : public — flush inline notes (if any) then open next scene within an act.
%          Optional argument accepted but ignored (historical).
\newcommand{\scene}[1][]{%
  \ifdefined\AVECSCENES  % no-op when actes only
  \ifdefined\AVECNOTES
    \ifdefined\ENDOFBOOKNOTES\else
      \WFclear% flush pending wrapfigs before page breaks
      \cleardoublepage
      % Reset after cleardoublepage: output routine may re-corrupt these via afterpage/floats
      \global\linewidth\textwidth
      \global\@totalleftmargin\z@
      {\centering\scenotitlefont Notes\par}%
      \vspace{\onelineskip}%
      \begingroup\parindent=0pt
      \theendnotes
      \endgroup
    \fi
  \fi
  \scene@init%
  \fi
}
% \scene@notes : reset endnote counter and add scene header to notes stream.
\newcommand{\scene@notes}[1]{%
  \ifdefined\AVECNOTES
    \setcounter{endnote}{0}%
    \ifdefined\ENDOFBOOKNOTES
      \addtoendnotes{\protect\par\protect\medskip\protect\noindent\protect\scene@noteheader{#1}\protect\par\protect\smallskip}%
    \fi
  \fi
}
\makeatother

% ─── PAGE GEOMETRY (155×235 mm – Medium 1) ───────────────────────────────────
\setstocksize{235mm}{155mm}
\settrimmedsize{235mm}{155mm}{*}
\settrims{0mm}{0mm}
\setlrmarginsandblock{23mm}{18mm}{*} % inner/outer
\setulmarginsandblock{25mm}{28mm}{*} % top/bottom
\setheadfoot{11mm}{13mm}
\setheaderspaces{*}{7mm}{*}
\checkandfixthelayout

% ─── TYPOGRAPHY ───────────────────────────────────────────────────────────────
\raggedbottom
\renewcommand{\baselinestretch}{1.2}
\setlength{\parindent}{4mm}
\setlength{\parskip}{0pt}
\setlength{\afterchapskip}{0.4\onelineskip}

% ─── DIALOGUE ─────────────────────────────────────────────────────────────────
\newenvironment{dialogue}
  {\par\vspace{0.5\baselineskip}\setlength{\parindent}{0pt}}
  {\par\vspace{0.5\baselineskip}}


% ─── TABLE DES MATIÈRES ───────────────────────────────────────────────────────
\setcounter{tocdepth}{1}
\renewcommand{\chapternumberline}[1]{}           % suppress arabic number, keep Roman title only
\renewcommand{\cftchapterdotsep}{\cftdotsep}     % dot leaders
\renewcommand{\cftdot}{{\color{black!35}.}}      % lighter dots throughout TOC
\renewcommand{\cftchapterfont}{\normalfont\scshape}     % chapter entry: small caps, not bold
\renewcommand{\cftchapterpagefont}{\normalfont\scshape} % page number: small caps, not bold
% Section entries = scenes (indented, normal font)
\setlength{\cftsectionindent}{2em}
\setlength{\cftsectionnumwidth}{0em}
\renewcommand{\cftsectionfont}{\normalfont}
\renewcommand{\cftsectionpagefont}{\normalfont}
\renewcommand{\cftsectiondotsep}{\cftdotsep}
% TOC title — gallimard chapter style suppresses \printchaptertitle, so we override here
\renewcommand{\printtoctitle}[1]{{\centering\normalfont\scshape\large #1\par}}
\renewcommand{\aftertoctitle}{\par\vspace{\onelineskip}}
% Table des extraits musicaux
\newlistof{listofscores}{mus}{Table des extraits musicaux}
\newlistentry{score}{mus}{0}
\setlength{\cftscoreindent}{0pt}
\newcommand{\mylistofscores}{%
  \cleardoublepage
  \phantomsection\addcontentsline{toc}{chapter}{Table des extraits musicaux}%
  \thispagestyle{empty}%
  \vspace*{4\onelineskip}%
  {\centering\normalfont\scshape\large Table des extraits musicaux\par}%
  \vspace{2\onelineskip}%
  \begingroup
  \small\renewcommand{\baselinestretch}{1.0}\selectfont
  \listofscores*%
  \endgroup}

% \mytableofcontents : cleardoublepage + suppress self-reference entry
\newcommand{\mytableofcontents}{%
  \cleardoublepage
  \begingroup
  \setlength{\beforechapskip}{0pt}%
  \small\renewcommand{\baselinestretch}{1.0}\selectfont
  \addtocontents{toc}{\protect\setcounter{tocdepth}{-1}}%
  \tableofcontents
  \addtocontents{toc}{\protect\setcounter{tocdepth}{1}}%
  \endgroup
}

% ─── ACTE / SCENE COUNTERS ────────────────────────────────────────────────────
\newcounter{actenum}   % incremented by \acte@init
\newcounter{scenenum}  % global, incremented by \scene@init
% shared font for scene-opener number and notes title
\newcommand{\scenotitlefont}{\normalfont\large}

% ─── CHAPTER STYLE ────────────────────────────────────────────────────────────
\makeatletter
\makechapterstyle{gallimard}{%
  \renewcommand{\chapnamefont}{\normalfont\scshape\centering}
  \renewcommand{\chapnumfont}{\scenotitlefont\centering}
  \renewcommand{\chaptitlefont}{\normalfont\scshape\centering}
  \renewcommand{\printchaptername}{}
  \renewcommand{\chapternamenum}{}
  \renewcommand{\printchapternum}{%
    \ifdefined\AVECSCENES
      \chapnumfont\arabic{scenenum}%
      \markboth{\Roman{actenum}}{\arabic{scenenum}}%
    \else
      \markboth{\Roman{actenum}}{}%
    \fi}
  \renewcommand{\afterchapternum}{%
    \ifdefined\AVECSCENES\par\vskip 0.4\onelineskip\fi}
  \renewcommand{\printchaptertitle}[1]{}
}
\makeatother
\chapterstyle{gallimard}

% ─── SCENE@INIT / ACTE@INIT / ACTE@TERM ──────────────────────────────────────
\makeatletter
% \scene@openLaTeXchapter : open a LaTeX chapter hidden from TOC, then init scene notes.
\newcommand{\scene@openLaTeXchapter}[1]{%
  \addtocontents{toc}{\protect\setcounter{tocdepth}{-1}}%
  \chapter{}%
  \addtocontents{toc}{\protect\setcounter{tocdepth}{1}}%
  \scene@notes{#1}%
}
% \scene@init : increment scene counter, compute notes label, open LaTeX chapter.
\newcommand{\scene@init}{%
  \ifdefined\AVECSCENES
    \stepcounter{scenenum}%
    \ifdefined\ENDOFBOOKNOTES
      \def\scene@noteslabel{Acte~\Roman{actenum}~·~Scène~\arabic{scenenum}}%
    \else
      \def\scene@noteslabel{Acte~\Roman{actenum}~·~Scène~\arabic{scenenum}}%
    \fi
  \else
    \ifdefined\ENDOFBOOKNOTES
      \def\scene@noteslabel{Acte~\Roman{actenum}}%
    \else
      \def\scene@noteslabel{\Roman{actenum}}%
    \fi
  \fi
  \scene@openLaTeXchapter{\scene@noteslabel}%
  \ifdefined\AVECSCENES
    \addcontentsline{toc}{section}{Scène~\arabic{scenenum}}%
  \fi
}
% \acte@term : flush notes from last scene of this act (no-op when \ENDOFBOOKNOTES or \AVECSCENES absent).
\newcommand{\acte@term}{%
  \ifdefined\AVECNOTES
    \ifdefined\ENDOFBOOKNOTES\else
      \ifdefined\AVECSCENES
        \WFclear% flush pending wrapfigs before page breaks
        \cleardoublepage
        % Reset after cleardoublepage: output routine may re-corrupt these via afterpage/floats
        \global\linewidth\textwidth
        \global\@totalleftmargin\z@
        {\centering\scenotitlefont Notes\par}%
        \vspace{\onelineskip}%
        \begingroup\parindent=0pt
        \theendnotes
        \endgroup
      \fi
    \fi
  \fi
}
% \acte@init : print act heading page and open first scene.
\newcommand{\acte@init}{%
  \stepcounter{actenum}%
  \cleartorecto%
  \thispagestyle{empty}%
  \vspace*{\fill}%
  {\centering\normalfont\scshape\large Acte~\Roman{actenum}\par}%
  \vspace*{\fill}%
  \phantomsection\addcontentsline{toc}{chapter}{Acte~\Roman{actenum}}%
  \clearpage%
  \scene@init%
}
% \acte : public — full act: heading page + content + flush notes.
\newcommand{\acte}[1]{\acte@init#1\acte@term}
\makeatother

% ─── RUNNING HEADERS : verso = acte centré pc, recto = scène alignée à droite ──
\makepagestyle{gallimard}
\makeevenhead{gallimard}{}{\footnotesize\scshape\leftmark}{}
\makeoddhead{gallimard}{}{}{\footnotesize\rightmark}
\makeevenfoot{gallimard}{}{\thepage}{}
\makeoddfoot{gallimard}{}{\thepage}{}
\makepsmarks{gallimard}{%
  \nouppercaseheads
}
\pagestyle{gallimard}
\aliaspagestyle{chapter}{gallimard}

% ─── EDITORIAL NOTE (Avertissement, Postfaces) ────────────────────────────────
% 10.95pt (~small), interligne ×1.1, sans alinéa, titre en petites capitales.
\newenvironment{editorialnote}[1]{%
  \begingroup
  \leftskip=0pt\relax\rightskip=0pt\relax
  \renewcommand{\baselinestretch}{1.1}%
  \small
  \setlength{\parindent}{0pt}%
  \setlength{\leftmargini}{0pt}%
  \setlength{\parskip}{\onelineskip}%
  \ifx&#1&\else
    {\centering\normalfont\scshape\large #1\par}%
    \vspace{\onelineskip}%
  \fi
  \noindent\ignorespaces
}{%
  \par\endgroup
}

%   \begin{document}
%   ...
% ─────────────────────────────────────────────────────────────────────────────

% ─── PDF VERSION ──────────────────────────────────────────────────────────────
\ifdefined\pdfvariable\pdfvariable minorversion 7\fi

% ─── PACKAGES ─────────────────────────────────────────────────────────────────
\usepackage{fontspec}
% ────── EB Garamond
%%% \setmainfont{EB Garamond}
% ────── Cormorant Garamond
%%% \setmainfont[
%%%   UprightFont    = *-Regular,
%%%   ItalicFont     = *-Italic,
%%%   BoldFont       = *-SemiBold,
%%%   BoldItalicFont = *-SemiBoldItalic
%%% ]{Cormorant Garamond}
% ────── EB Garamond (main) + Adobe Garamond Pro (bold only)
\setmainfont{EB Garamond}[
  BoldFont       = AGaramondPro-Semibold,
  BoldItalicFont = AGaramondPro-SemiboldItalic,
]

\newfontfamily{\arabicfont}{Geeza Pro}
\newcommand{\arab}[1]{{\arabicfont\textdir TRT #1}}
\newfontfamily{\cjkfont}{Hiragino Mincho ProN}
\newcommand{\cjk}[1]{{\cjkfont #1}}
\newcommand{\shouting}[1]{\textbf{\textsc{#1}}}
\usepackage[russian,french]{babel}
\usepackage[program=/usr/local/bin/lilypond,includepaths={music}]{lyluatex}
\usepackage{microtype}
\usepackage{fmtcount}
\usepackage{catchfile}
\usepackage{graphicx}
\usepackage{xcolor}
\usepackage{amsmath}
\usepackage{xparse}
\usepackage{tikz}
\usepackage{afterpage}
\usepackage{needspace}
\usepackage{pdfpages}
\usepackage[absolute,overlay]{textpos}
\usepackage{tabularx}
\usepackage{booktabs}
\usepackage{calc}
\usepackage{wrapfig}
\usepackage{tcolorbox}
\usepackage{ragged2e}
% Fix: \WF@finale restores \linewidth and \@totalleftmargin only when
% \parshape==\WF@fudgeparshape; the "forced to float" collision path bypasses
% that guard, leaving both corrupted globally.  Patch \WF@finale to force-reset
% both after every call, regardless of which path was taken.
\makeatletter
\let\WF@@orig@finale\WF@finale
\def\WF@finale{\WF@@orig@finale\global\linewidth\textwidth\global\@totalleftmargin\z@}
\makeatother
% Safety net: clear any residual wrapfig state at chapter boundaries.
\pretocmd{\chapter}{\WFclear}{}{}
\usepackage{hyperref}
\usepackage{bookmark}
\hypersetup{pdfpagelayout=TwoPageRight, hidelinks}
\usetikzlibrary{calc,positioning}


% ─── VERSION / GIT ────────────────────────────────────────────────────────────
\newcommand{\docversion}{0.5.0}
\newcommand{\docversionordinal}{Deuxième}

\IfFileExists{tete_de_veau_ravigote.gitinfo}%
  {\CatchFileDef{\gitcommit}{tete_de_veau_ravigote.gitinfo}{\endlinechar=-1}}%
  {\def\gitcommit{}}

% ─── GRAYSCALE IMAGE FILTER ───────────────────────────────────────────────────
\graphicspath{{iconography/}}
\newcommand{\figcaptionname}{\footnotesize}
\newcommand{\figcaptiondetail}{\scriptsize}
\newcommand{\figpagetitlename}{\small\scshape}
\newcommand{\figpagetitledetail}{\footnotesize}
\newcommand{\bwimage}[2][]{%
  \begin{tikzpicture}
    \node[inner sep=0, outer sep=0pt, fill=white] (img) {\includegraphics[#1]{#2}};
    %%% begin{scope}[blend mode=saturation]
    %%%  \fill[black] (img.south west) rectangle (img.north east);
    %%% end{scope}
  \end{tikzpicture}%
}

% ─── ESPACEMENT FLOTTANTS ─────────────────────────────────────────────────────
\setlength{\textfloatsep}{8pt}   % espace entre flottant top/bottom et le texte (défaut ~20pt)

% ─── ICONOGRAPHIE ─────────────────────────────────────────────────────────────
% All commands below are no-ops unless \AVECILLUSTRATIONS is defined.
% Image paths are relative to the graphicspath (iconography/).
\ifdefined\AVECILLUSTRATIONS
  \newlength{\iconofillheight}% scratch register for [p] page-fill placement
  % \iconographieimg[scale]{image}{caption}
  % Single image centred on its own dedicated page (via \afterpage).
  % [scale] height as fraction of \textheight (default 0.75); {caption} may be empty.
  \newcommand{\iconographieimg}[3][0.75]{%
    \afterpage{%
      \thispagestyle{gallimard}%
      \vspace*{\fill}%
      \begin{center}%
        \bwimage[height=#1\textheight,width=\textwidth,keepaspectratio]{#2}\\[0.6em]%
        {\figcaptionname\itshape #3}%
      \end{center}%
      \vspace*{\fill}%
      \newpage
    }%
  }
  % \iconographietex{figures/foo_fig.tex}
  % Inserts one hand-crafted figure page from an external .tex file (via \afterpage).
  % The file is responsible for its own layout — see figures/*_fig.tex for examples.
  \newcommand{\iconographietex}[1]{%
    \afterpage{%
      \thispagestyle{gallimard}%
      \input{#1}%
      \clearpage
    }%
  }
  % \iconographietexpair{figures/foo_fig.tex}{figures/bar_fig.tex}
  % Like \iconographietex but inserts two consecutive figure pages.
  \newcommand{\iconographietexpair}[2]{%
    \afterpage{%
      \thispagestyle{gallimard}%
      \input{#1}%
      \input{#2}%
      \clearpage
    }%
  }
  % \iconographiewrapfig{pos}{width}{\bwimage[width=\linewidth]{path}}{caption}
  % Image wrapped by surrounding text (wrapfigure).
  % {pos} = l or r;  {width} = e.g. 0.35\textwidth;  {caption} may be empty {}.
  \newcommand{\iconographiewrapfig}[4]{%
    \setlength{\intextsep}{0pt}%
    \setlength{\columnsep}{8pt}%
    \begin{wrapfigure}{#1}{#2}%
      \centering
      #3%
      \ifx&#4&\else\\[0.2em]{\figcaptionname\itshape #4}\\[0.4\baselineskip]\fi%
    \end{wrapfigure}%
  }
  % \iconographieinlineblock[placement][lines]{\bwimage[width=W]{path}}{caption}
  % Centred image block in the text flow. {caption} may be empty {}.
  % [placement] empty or omitted → inline: caller must put \par before and \noindent after.
  % [placement] = t or b → LaTeX float (top/bottom of page); no \par/\noindent needed.
  % [placement] = p → fill remaining page height inline (image scaled to fit, no float).
  % [lines] number of caption lines (default 1); used by [p] to reserve the right height.
  % Multiple inline calls in the same paragraph are supported.
  \NewDocumentCommand{\iconographieinlineblock}{O{} O{1} m m}{%
    % #1=placement  #2=caption lines (for [p])  #3=image  #4=caption
    \ifx&#1&%
      % no placement → inline
      \vspace{0.5\onelineskip}%
      \centerline{#3}%
      \ifx&#4&\else\par\centerline{\figcaptionname\itshape#4}\vspace{0.5\baselineskip}\fi%
      \vspace{-0.5\onelineskip}%
    \else
      \ifthenelse{\equal{#1}{p}}{%
        % p → scale image to fill remaining page height
        \par
        \begingroup
          \setlength{\iconofillheight}{\dimexpr\pagegoal-\pagetotal\relax}%
          \addtolength{\iconofillheight}{-2\baselineskip}%
          \ifx&#4&\else\addtolength{\iconofillheight}{-#2\baselineskip-0.5\baselineskip}\fi%
          \ifdim\iconofillheight>0pt
            \vspace{2\baselineskip}%
            \centerline{\resizebox*{!}{\iconofillheight}{#3}}%
            \ifx&#4&\else\par{\centering\figcaptionname\itshape#4\par}\fi%
          \fi
        \endgroup
      }{%
        % placement given → float
        \begin{figure}[#1]
          \centering#3%
          \ifx&#4&\else\par{\figcaptionname\itshape#4}\fi%
        \end{figure}%
      }%
    \fi
  }
  % \iconographiedouble[placement][height][margin]{\bwimage{img1}}{\bwimage{img2}}[caption]
  % Two images side by side at equal height.
  % [placement] = b (default, push to bottom) or t (push to top)
  % [height] common image height (default 5cm); [margin] symmetric indent (default 0pt)
  % [caption] shared caption, omit for none.
  \NewDocumentCommand{\iconographiedouble}{O{b} O{5cm} O{0pt} m m O{}}{%
    \ifthenelse{\equal{#1}{b}}{\vspace*{\fill}}{\vspace*{0pt}}%
    \vspace{0.5\onelineskip}%
    {\setlength{\leftskip}{#3}\setlength{\rightskip}{#3}%
    \noindent\bwimage[height=#2]{#4}\hfill\bwimage[height=#2]{#5}\par}%
    \ifx&#6&%
      \vspace{0.5\onelineskip}%
    \else%
      \vspace{0.2\onelineskip}%
      {\centering\figcaptionname\itshape#6\par}%
      \vspace{0.3\onelineskip}%
    \fi%
    \ifthenelse{\equal{#1}{t}}{\vspace*{\fill}}{}%
  }
  % \iconographietriple[placement][height][margin]{\bwimage{img1}}{\bwimage{img2}}{\bwimage{img3}}[caption]
  % Three images side by side at equal height. Same options as \iconographiedouble.
  \NewDocumentCommand{\iconographietriple}{O{b} O{5cm} O{0pt} m m m O{}}{%
    \ifthenelse{\equal{#1}{b}}{\vspace*{\fill}}{\vspace*{0pt}}%
    \vspace{0.5\onelineskip}%
    {\setlength{\leftskip}{#3}\setlength{\rightskip}{#3}%
    \noindent\bwimage[height=#2]{#4}\hfill\bwimage[height=#2]{#5}\hfill\bwimage[height=#2]{#6}\par}%
    \ifx&#7&%
      \vspace{0.5\onelineskip}%
    \else%
      \vspace{0.2\onelineskip}%
      {\centering\figcaptionname\itshape#7\par}%
      \vspace{0.3\onelineskip}%
    \fi%
    \ifthenelse{\equal{#1}{t}}{\vspace*{\fill}}{}%
  }
\else
  \newcommand{\iconographieimg}[3][0.75]{}
  \newcommand{\iconographietex}[1]{}
  \newcommand{\iconographietexpair}[2]{}
  \newcommand{\iconographiewrapfig}[4]{}
  \NewDocumentCommand{\iconographieinlineblock}{O{} O{1} m m}{}
  \NewDocumentCommand{\iconographiedouble}{O{b} O{5cm} O{0pt} m m O{}}{}
  \NewDocumentCommand{\iconographietriple}{O{b} O{5cm} O{0pt} m m m O{}}{}
\fi

% ─── EXTRAITS MUSICAUX ────────────────────────────────────────────────────────
% Compiler avec : lualatex "\def\AVECMUSIQUE{}\input{...}"
\ifdefined\AVECMUSIQUE
  % Extrait musical inline : [#1]=espacement (défaut 0.66), #2=fichier dans music/, #3=légende
  \newcommand{\score}[3][0.66]{%
    \addcontentsline{mus}{score}{#3}%
    \noindent\begin{minipage}{\textwidth}%
    \vspace{#1\onelineskip}%
    \lilypondfile[line-width=\textwidth]{music/#2}%
    \begin{center}%
      \vspace{-#1\onelineskip}%
      {\centering\figcaptionname\itshape#3}%
      \vspace{#1\onelineskip}%
    \end{center}%
    \end{minipage}}
\else
  \newcommand{\score}[3][1]{}
\fi

% ─── NOTES FINALES ────────────────────────────────────────────────────────────
% Compiler avec : lualatex "\def\AVECNOTES{}\input{...}"
% \ENDOFBOOKNOTES : notes regroupées en fin de livre (avec en-tête de scène) ;
%                   absent = notes à la fin de chaque scène (sans en-tête).
\usepackage{endnotes}
\ifdefined\AVECNOTES
  \newcommand{\nf}[1]{\endnote{#1}}
\else
  \newcommand{\nf}[1]{}
\fi
% Supprimer le titre automatique du package (on gère le nôtre)
\renewcommand{\notesname}{}
% Style des notes : même fonte, corps légèrement réduit
\renewcommand{\enotesize}{\small}
% Format des notes en fin de livre : sans retrait, numéro en corps normal
\makeatletter
\renewcommand\enoteformat{%
  \rightskip\z@ \leftskip\z@ \parindent=0pt
  \leavevmode\theenmark.\enspace}
\makeatother
% Fix PDF ActualText for œ (oe-ligature) copy-paste in \source URLs
\usepackage{accsupp}
\newcommand{\oechar}{\char"153\relax}
{\catcode"153=\active\gdefœ{\BeginAccSupp{method=hex,unicode,ActualText=0153}\oechar\EndAccSupp{}}}
% Source dans les notes : saut de ligne avant le lien Wikipedia
\newcommand{\source}[1]{\newline\textit{Source~:}~{\upshape\addfontfeatures{RawFeature=-tlig;-calt}\def\_{\textunderscore\allowbreak}\def\%{\char"25\relax\allowbreak}\catcode"153=\active\scantokens{#1\empty}}}
% Image seule dans une note (centrée sous le texte)
\newcommand{\nfimg}[2]{\newline\includegraphics[width=#1\linewidth]{#2}}
% Image côte à côte avec le texte : #1=texte, #2=frac image, #3=path, #4=url source
\newcommand{\nfimgblock}[4]{%
  \begin{minipage}[t]{0.45\linewidth}\small#1\end{minipage}%
  \hfill%
  \begin{minipage}[t]{#2\linewidth}\includegraphics[width=\linewidth]{#3}\end{minipage}%
  \source{#4}}
% Filet ornemental (séparateur de scène intra-chapitre)
\newcommand{\ornamentalrule}{%
  \vspace{\onelineskip}%
  {\centering\rule{0.4\textwidth}{0.4pt}\par}%
  \vspace{\onelineskip}%
}
% Font for scene header in notes block
\makeatletter
\ifdefined\ENDOFBOOKNOTES
  \newcommand{\scene@noteheader}[1]{{\normalfont\large #1}}
\else
  \newcommand{\scene@noteheader}[1]{{\normalfont\scshape #1}}
\fi
% \scene : public — flush inline notes (if any) then open next scene within an act.
%          Optional argument accepted but ignored (historical).
\newcommand{\scene}[1][]{%
  \ifdefined\AVECSCENES  % no-op when actes only
  \ifdefined\AVECNOTES
    \ifdefined\ENDOFBOOKNOTES\else
      \WFclear% flush pending wrapfigs before page breaks
      \cleardoublepage
      % Reset after cleardoublepage: output routine may re-corrupt these via afterpage/floats
      \global\linewidth\textwidth
      \global\@totalleftmargin\z@
      {\centering\scenotitlefont Notes\par}%
      \vspace{\onelineskip}%
      \begingroup\parindent=0pt
      \theendnotes
      \endgroup
    \fi
  \fi
  \scene@init%
  \fi
}
% \scene@notes : reset endnote counter and add scene header to notes stream.
\newcommand{\scene@notes}[1]{%
  \ifdefined\AVECNOTES
    \setcounter{endnote}{0}%
    \ifdefined\ENDOFBOOKNOTES
      \addtoendnotes{\protect\par\protect\medskip\protect\noindent\protect\scene@noteheader{#1}\protect\par\protect\smallskip}%
    \fi
  \fi
}
\makeatother

% ─── PAGE GEOMETRY (155×235 mm – Medium 1) ───────────────────────────────────
\setstocksize{235mm}{155mm}
\settrimmedsize{235mm}{155mm}{*}
\settrims{0mm}{0mm}
\setlrmarginsandblock{23mm}{18mm}{*} % inner/outer
\setulmarginsandblock{25mm}{28mm}{*} % top/bottom
\setheadfoot{11mm}{13mm}
\setheaderspaces{*}{7mm}{*}
\checkandfixthelayout

% ─── TYPOGRAPHY ───────────────────────────────────────────────────────────────
\raggedbottom
\renewcommand{\baselinestretch}{1.2}
\setlength{\parindent}{4mm}
\setlength{\parskip}{0pt}
\setlength{\afterchapskip}{0.4\onelineskip}

% ─── DIALOGUE ─────────────────────────────────────────────────────────────────
\newenvironment{dialogue}
  {\par\vspace{0.5\baselineskip}\setlength{\parindent}{0pt}}
  {\par\vspace{0.5\baselineskip}}


% ─── TABLE DES MATIÈRES ───────────────────────────────────────────────────────
\setcounter{tocdepth}{1}
\renewcommand{\chapternumberline}[1]{}           % suppress arabic number, keep Roman title only
\renewcommand{\cftchapterdotsep}{\cftdotsep}     % dot leaders
\renewcommand{\cftdot}{{\color{black!35}.}}      % lighter dots throughout TOC
\renewcommand{\cftchapterfont}{\normalfont\scshape}     % chapter entry: small caps, not bold
\renewcommand{\cftchapterpagefont}{\normalfont\scshape} % page number: small caps, not bold
% Section entries = scenes (indented, normal font)
\setlength{\cftsectionindent}{2em}
\setlength{\cftsectionnumwidth}{0em}
\renewcommand{\cftsectionfont}{\normalfont}
\renewcommand{\cftsectionpagefont}{\normalfont}
\renewcommand{\cftsectiondotsep}{\cftdotsep}
% TOC title — gallimard chapter style suppresses \printchaptertitle, so we override here
\renewcommand{\printtoctitle}[1]{{\centering\normalfont\scshape\large #1\par}}
\renewcommand{\aftertoctitle}{\par\vspace{\onelineskip}}
% Table des extraits musicaux
\newlistof{listofscores}{mus}{Table des extraits musicaux}
\newlistentry{score}{mus}{0}
\setlength{\cftscoreindent}{0pt}
\newcommand{\mylistofscores}{%
  \cleardoublepage
  \phantomsection\addcontentsline{toc}{chapter}{Table des extraits musicaux}%
  \thispagestyle{empty}%
  \vspace*{4\onelineskip}%
  {\centering\normalfont\scshape\large Table des extraits musicaux\par}%
  \vspace{2\onelineskip}%
  \begingroup
  \small\renewcommand{\baselinestretch}{1.0}\selectfont
  \listofscores*%
  \endgroup}

% \mytableofcontents : cleardoublepage + suppress self-reference entry
\newcommand{\mytableofcontents}{%
  \cleardoublepage
  \begingroup
  \setlength{\beforechapskip}{0pt}%
  \small\renewcommand{\baselinestretch}{1.0}\selectfont
  \addtocontents{toc}{\protect\setcounter{tocdepth}{-1}}%
  \tableofcontents
  \addtocontents{toc}{\protect\setcounter{tocdepth}{1}}%
  \endgroup
}

% ─── ACTE / SCENE COUNTERS ────────────────────────────────────────────────────
\newcounter{actenum}   % incremented by \acte@init
\newcounter{scenenum}  % global, incremented by \scene@init
% shared font for scene-opener number and notes title
\newcommand{\scenotitlefont}{\normalfont\large}

% ─── CHAPTER STYLE ────────────────────────────────────────────────────────────
\makeatletter
\makechapterstyle{gallimard}{%
  \renewcommand{\chapnamefont}{\normalfont\scshape\centering}
  \renewcommand{\chapnumfont}{\scenotitlefont\centering}
  \renewcommand{\chaptitlefont}{\normalfont\scshape\centering}
  \renewcommand{\printchaptername}{}
  \renewcommand{\chapternamenum}{}
  \renewcommand{\printchapternum}{%
    \ifdefined\AVECSCENES
      \chapnumfont\arabic{scenenum}%
      \markboth{\Roman{actenum}}{\arabic{scenenum}}%
    \else
      \markboth{\Roman{actenum}}{}%
    \fi}
  \renewcommand{\afterchapternum}{%
    \ifdefined\AVECSCENES\par\vskip 0.4\onelineskip\fi}
  \renewcommand{\printchaptertitle}[1]{}
}
\makeatother
\chapterstyle{gallimard}

% ─── SCENE@INIT / ACTE@INIT / ACTE@TERM ──────────────────────────────────────
\makeatletter
% \scene@openLaTeXchapter : open a LaTeX chapter hidden from TOC, then init scene notes.
\newcommand{\scene@openLaTeXchapter}[1]{%
  \addtocontents{toc}{\protect\setcounter{tocdepth}{-1}}%
  \chapter{}%
  \addtocontents{toc}{\protect\setcounter{tocdepth}{1}}%
  \scene@notes{#1}%
}
% \scene@init : increment scene counter, compute notes label, open LaTeX chapter.
\newcommand{\scene@init}{%
  \ifdefined\AVECSCENES
    \stepcounter{scenenum}%
    \ifdefined\ENDOFBOOKNOTES
      \def\scene@noteslabel{Acte~\Roman{actenum}~·~Scène~\arabic{scenenum}}%
    \else
      \def\scene@noteslabel{Acte~\Roman{actenum}~·~Scène~\arabic{scenenum}}%
    \fi
  \else
    \ifdefined\ENDOFBOOKNOTES
      \def\scene@noteslabel{Acte~\Roman{actenum}}%
    \else
      \def\scene@noteslabel{\Roman{actenum}}%
    \fi
  \fi
  \scene@openLaTeXchapter{\scene@noteslabel}%
  \ifdefined\AVECSCENES
    \addcontentsline{toc}{section}{Scène~\arabic{scenenum}}%
  \fi
}
% \acte@term : flush notes from last scene of this act (no-op when \ENDOFBOOKNOTES or \AVECSCENES absent).
\newcommand{\acte@term}{%
  \ifdefined\AVECNOTES
    \ifdefined\ENDOFBOOKNOTES\else
      \ifdefined\AVECSCENES
        \WFclear% flush pending wrapfigs before page breaks
        \cleardoublepage
        % Reset after cleardoublepage: output routine may re-corrupt these via afterpage/floats
        \global\linewidth\textwidth
        \global\@totalleftmargin\z@
        {\centering\scenotitlefont Notes\par}%
        \vspace{\onelineskip}%
        \begingroup\parindent=0pt
        \theendnotes
        \endgroup
      \fi
    \fi
  \fi
}
% \acte@init : print act heading page and open first scene.
\newcommand{\acte@init}{%
  \stepcounter{actenum}%
  \cleartorecto%
  \thispagestyle{empty}%
  \vspace*{\fill}%
  {\centering\normalfont\scshape\large Acte~\Roman{actenum}\par}%
  \vspace*{\fill}%
  \phantomsection\addcontentsline{toc}{chapter}{Acte~\Roman{actenum}}%
  \clearpage%
  \scene@init%
}
% \acte : public — full act: heading page + content + flush notes.
\newcommand{\acte}[1]{\acte@init#1\acte@term}
\makeatother

% ─── RUNNING HEADERS : verso = acte centré pc, recto = scène alignée à droite ──
\makepagestyle{gallimard}
\makeevenhead{gallimard}{}{\footnotesize\scshape\leftmark}{}
\makeoddhead{gallimard}{}{}{\footnotesize\rightmark}
\makeevenfoot{gallimard}{}{\thepage}{}
\makeoddfoot{gallimard}{}{\thepage}{}
\makepsmarks{gallimard}{%
  \nouppercaseheads
}
\pagestyle{gallimard}
\aliaspagestyle{chapter}{gallimard}

% ─── EDITORIAL NOTE (Avertissement, Postfaces) ────────────────────────────────
% 10.95pt (~small), interligne ×1.1, sans alinéa, titre en petites capitales.
\newenvironment{editorialnote}[1]{%
  \begingroup
  \leftskip=0pt\relax\rightskip=0pt\relax
  \renewcommand{\baselinestretch}{1.1}%
  \small
  \setlength{\parindent}{0pt}%
  \setlength{\leftmargini}{0pt}%
  \setlength{\parskip}{\onelineskip}%
  \ifx&#1&\else
    {\centering\normalfont\scshape\large #1\par}%
    \vspace{\onelineskip}%
  \fi
  \noindent\ignorespaces
}{%
  \par\endgroup
}

%   \begin{document}
%   ...
% ─────────────────────────────────────────────────────────────────────────────

% ─── PDF VERSION ──────────────────────────────────────────────────────────────
\ifdefined\pdfvariable\pdfvariable minorversion 7\fi

% ─── PACKAGES ─────────────────────────────────────────────────────────────────
\usepackage{fontspec}
% ────── EB Garamond
%%% \setmainfont{EB Garamond}
% ────── Cormorant Garamond
%%% \setmainfont[
%%%   UprightFont    = *-Regular,
%%%   ItalicFont     = *-Italic,
%%%   BoldFont       = *-SemiBold,
%%%   BoldItalicFont = *-SemiBoldItalic
%%% ]{Cormorant Garamond}
% ────── EB Garamond (main) + Adobe Garamond Pro (bold only)
\setmainfont{EB Garamond}[
  BoldFont       = AGaramondPro-Semibold,
  BoldItalicFont = AGaramondPro-SemiboldItalic,
]

\newfontfamily{\arabicfont}{Geeza Pro}
\newcommand{\arab}[1]{{\arabicfont\textdir TRT #1}}
\newfontfamily{\cjkfont}{Hiragino Mincho ProN}
\newcommand{\cjk}[1]{{\cjkfont #1}}
\newcommand{\shouting}[1]{\textbf{\textsc{#1}}}
\usepackage[russian,french]{babel}
\usepackage[program=/usr/local/bin/lilypond,includepaths={music}]{lyluatex}
\usepackage{microtype}
\usepackage{fmtcount}
\usepackage{catchfile}
\usepackage{graphicx}
\usepackage{xcolor}
\usepackage{amsmath}
\usepackage{xparse}
\usepackage{tikz}
\usepackage{afterpage}
\usepackage{needspace}
\usepackage{pdfpages}
\usepackage[absolute,overlay]{textpos}
\usepackage{tabularx}
\usepackage{booktabs}
\usepackage{calc}
\usepackage{wrapfig}
\usepackage{tcolorbox}
\usepackage{ragged2e}
% Fix: \WF@finale restores \linewidth and \@totalleftmargin only when
% \parshape==\WF@fudgeparshape; the "forced to float" collision path bypasses
% that guard, leaving both corrupted globally.  Patch \WF@finale to force-reset
% both after every call, regardless of which path was taken.
\makeatletter
\let\WF@@orig@finale\WF@finale
\def\WF@finale{\WF@@orig@finale\global\linewidth\textwidth\global\@totalleftmargin\z@}
\makeatother
% Safety net: clear any residual wrapfig state at chapter boundaries.
\pretocmd{\chapter}{\WFclear}{}{}
\usepackage{hyperref}
\usepackage{bookmark}
\hypersetup{pdfpagelayout=TwoPageRight, hidelinks}
\usetikzlibrary{calc,positioning}


% ─── VERSION / GIT ────────────────────────────────────────────────────────────
\newcommand{\docversion}{0.5.0}
\newcommand{\docversionordinal}{Deuxième}

\IfFileExists{tete_de_veau_ravigote.gitinfo}%
  {\CatchFileDef{\gitcommit}{tete_de_veau_ravigote.gitinfo}{\endlinechar=-1}}%
  {\def\gitcommit{}}

% ─── GRAYSCALE IMAGE FILTER ───────────────────────────────────────────────────
\graphicspath{{iconography/}}
\newcommand{\figcaptionname}{\footnotesize}
\newcommand{\figcaptiondetail}{\scriptsize}
\newcommand{\figpagetitlename}{\small\scshape}
\newcommand{\figpagetitledetail}{\footnotesize}
\newcommand{\bwimage}[2][]{%
  \begin{tikzpicture}
    \node[inner sep=0, outer sep=0pt, fill=white] (img) {\includegraphics[#1]{#2}};
    %%% begin{scope}[blend mode=saturation]
    %%%  \fill[black] (img.south west) rectangle (img.north east);
    %%% end{scope}
  \end{tikzpicture}%
}

% ─── ESPACEMENT FLOTTANTS ─────────────────────────────────────────────────────
\setlength{\textfloatsep}{8pt}   % espace entre flottant top/bottom et le texte (défaut ~20pt)

% ─── ICONOGRAPHIE ─────────────────────────────────────────────────────────────
% All commands below are no-ops unless \AVECILLUSTRATIONS is defined.
% Image paths are relative to the graphicspath (iconography/).
\ifdefined\AVECILLUSTRATIONS
  \newlength{\iconofillheight}% scratch register for [p] page-fill placement
  % \iconographieimg[scale]{image}{caption}
  % Single image centred on its own dedicated page (via \afterpage).
  % [scale] height as fraction of \textheight (default 0.75); {caption} may be empty.
  \newcommand{\iconographieimg}[3][0.75]{%
    \afterpage{%
      \thispagestyle{gallimard}%
      \vspace*{\fill}%
      \begin{center}%
        \bwimage[height=#1\textheight,width=\textwidth,keepaspectratio]{#2}\\[0.6em]%
        {\figcaptionname\itshape #3}%
      \end{center}%
      \vspace*{\fill}%
      \newpage
    }%
  }
  % \iconographietex{figures/foo_fig.tex}
  % Inserts one hand-crafted figure page from an external .tex file (via \afterpage).
  % The file is responsible for its own layout — see figures/*_fig.tex for examples.
  \newcommand{\iconographietex}[1]{%
    \afterpage{%
      \thispagestyle{gallimard}%
      \input{#1}%
      \clearpage
    }%
  }
  % \iconographietexpair{figures/foo_fig.tex}{figures/bar_fig.tex}
  % Like \iconographietex but inserts two consecutive figure pages.
  \newcommand{\iconographietexpair}[2]{%
    \afterpage{%
      \thispagestyle{gallimard}%
      \input{#1}%
      \input{#2}%
      \clearpage
    }%
  }
  % \iconographiewrapfig{pos}{width}{\bwimage[width=\linewidth]{path}}{caption}
  % Image wrapped by surrounding text (wrapfigure).
  % {pos} = l or r;  {width} = e.g. 0.35\textwidth;  {caption} may be empty {}.
  \newcommand{\iconographiewrapfig}[4]{%
    \setlength{\intextsep}{0pt}%
    \setlength{\columnsep}{8pt}%
    \begin{wrapfigure}{#1}{#2}%
      \centering
      #3%
      \ifx&#4&\else\\[0.2em]{\figcaptionname\itshape #4}\\[0.4\baselineskip]\fi%
    \end{wrapfigure}%
  }
  % \iconographieinlineblock[placement][lines]{\bwimage[width=W]{path}}{caption}
  % Centred image block in the text flow. {caption} may be empty {}.
  % [placement] empty or omitted → inline: caller must put \par before and \noindent after.
  % [placement] = t or b → LaTeX float (top/bottom of page); no \par/\noindent needed.
  % [placement] = p → fill remaining page height inline (image scaled to fit, no float).
  % [lines] number of caption lines (default 1); used by [p] to reserve the right height.
  % Multiple inline calls in the same paragraph are supported.
  \NewDocumentCommand{\iconographieinlineblock}{O{} O{1} m m}{%
    % #1=placement  #2=caption lines (for [p])  #3=image  #4=caption
    \ifx&#1&%
      % no placement → inline
      \vspace{0.5\onelineskip}%
      \centerline{#3}%
      \ifx&#4&\else\par\centerline{\figcaptionname\itshape#4}\vspace{0.5\baselineskip}\fi%
      \vspace{-0.5\onelineskip}%
    \else
      \ifthenelse{\equal{#1}{p}}{%
        % p → scale image to fill remaining page height
        \par
        \begingroup
          \setlength{\iconofillheight}{\dimexpr\pagegoal-\pagetotal\relax}%
          \addtolength{\iconofillheight}{-2\baselineskip}%
          \ifx&#4&\else\addtolength{\iconofillheight}{-#2\baselineskip-0.5\baselineskip}\fi%
          \ifdim\iconofillheight>0pt
            \vspace{2\baselineskip}%
            \centerline{\resizebox*{!}{\iconofillheight}{#3}}%
            \ifx&#4&\else\par{\centering\figcaptionname\itshape#4\par}\fi%
          \fi
        \endgroup
      }{%
        % placement given → float
        \begin{figure}[#1]
          \centering#3%
          \ifx&#4&\else\par{\figcaptionname\itshape#4}\fi%
        \end{figure}%
      }%
    \fi
  }
  % \iconographiedouble[placement][height][margin]{\bwimage{img1}}{\bwimage{img2}}[caption]
  % Two images side by side at equal height.
  % [placement] = b (default, push to bottom) or t (push to top)
  % [height] common image height (default 5cm); [margin] symmetric indent (default 0pt)
  % [caption] shared caption, omit for none.
  \NewDocumentCommand{\iconographiedouble}{O{b} O{5cm} O{0pt} m m O{}}{%
    \ifthenelse{\equal{#1}{b}}{\vspace*{\fill}}{\vspace*{0pt}}%
    \vspace{0.5\onelineskip}%
    {\setlength{\leftskip}{#3}\setlength{\rightskip}{#3}%
    \noindent\bwimage[height=#2]{#4}\hfill\bwimage[height=#2]{#5}\par}%
    \ifx&#6&%
      \vspace{0.5\onelineskip}%
    \else%
      \vspace{0.2\onelineskip}%
      {\centering\figcaptionname\itshape#6\par}%
      \vspace{0.3\onelineskip}%
    \fi%
    \ifthenelse{\equal{#1}{t}}{\vspace*{\fill}}{}%
  }
  % \iconographietriple[placement][height][margin]{\bwimage{img1}}{\bwimage{img2}}{\bwimage{img3}}[caption]
  % Three images side by side at equal height. Same options as \iconographiedouble.
  \NewDocumentCommand{\iconographietriple}{O{b} O{5cm} O{0pt} m m m O{}}{%
    \ifthenelse{\equal{#1}{b}}{\vspace*{\fill}}{\vspace*{0pt}}%
    \vspace{0.5\onelineskip}%
    {\setlength{\leftskip}{#3}\setlength{\rightskip}{#3}%
    \noindent\bwimage[height=#2]{#4}\hfill\bwimage[height=#2]{#5}\hfill\bwimage[height=#2]{#6}\par}%
    \ifx&#7&%
      \vspace{0.5\onelineskip}%
    \else%
      \vspace{0.2\onelineskip}%
      {\centering\figcaptionname\itshape#7\par}%
      \vspace{0.3\onelineskip}%
    \fi%
    \ifthenelse{\equal{#1}{t}}{\vspace*{\fill}}{}%
  }
\else
  \newcommand{\iconographieimg}[3][0.75]{}
  \newcommand{\iconographietex}[1]{}
  \newcommand{\iconographietexpair}[2]{}
  \newcommand{\iconographiewrapfig}[4]{}
  \NewDocumentCommand{\iconographieinlineblock}{O{} O{1} m m}{}
  \NewDocumentCommand{\iconographiedouble}{O{b} O{5cm} O{0pt} m m O{}}{}
  \NewDocumentCommand{\iconographietriple}{O{b} O{5cm} O{0pt} m m m O{}}{}
\fi

% ─── EXTRAITS MUSICAUX ────────────────────────────────────────────────────────
% Compiler avec : lualatex "\def\AVECMUSIQUE{}\input{...}"
\ifdefined\AVECMUSIQUE
  % Extrait musical inline : [#1]=espacement (défaut 0.66), #2=fichier dans music/, #3=légende
  \newcommand{\score}[3][0.66]{%
    \addcontentsline{mus}{score}{#3}%
    \noindent\begin{minipage}{\textwidth}%
    \vspace{#1\onelineskip}%
    \lilypondfile[line-width=\textwidth]{music/#2}%
    \begin{center}%
      \vspace{-#1\onelineskip}%
      {\centering\figcaptionname\itshape#3}%
      \vspace{#1\onelineskip}%
    \end{center}%
    \end{minipage}}
\else
  \newcommand{\score}[3][1]{}
\fi

% ─── NOTES FINALES ────────────────────────────────────────────────────────────
% Compiler avec : lualatex "\def\AVECNOTES{}\input{...}"
% \ENDOFBOOKNOTES : notes regroupées en fin de livre (avec en-tête de scène) ;
%                   absent = notes à la fin de chaque scène (sans en-tête).
\usepackage{endnotes}
\ifdefined\AVECNOTES
  \newcommand{\nf}[1]{\endnote{#1}}
\else
  \newcommand{\nf}[1]{}
\fi
% Supprimer le titre automatique du package (on gère le nôtre)
\renewcommand{\notesname}{}
% Style des notes : même fonte, corps légèrement réduit
\renewcommand{\enotesize}{\small}
% Format des notes en fin de livre : sans retrait, numéro en corps normal
\makeatletter
\renewcommand\enoteformat{%
  \rightskip\z@ \leftskip\z@ \parindent=0pt
  \leavevmode\theenmark.\enspace}
\makeatother
% Fix PDF ActualText for œ (oe-ligature) copy-paste in \source URLs
\usepackage{accsupp}
\newcommand{\oechar}{\char"153\relax}
{\catcode"153=\active\gdefœ{\BeginAccSupp{method=hex,unicode,ActualText=0153}\oechar\EndAccSupp{}}}
% Source dans les notes : saut de ligne avant le lien Wikipedia
\newcommand{\source}[1]{\newline\textit{Source~:}~{\upshape\addfontfeatures{RawFeature=-tlig;-calt}\def\_{\textunderscore\allowbreak}\def\%{\char"25\relax\allowbreak}\catcode"153=\active\scantokens{#1\empty}}}
% Image seule dans une note (centrée sous le texte)
\newcommand{\nfimg}[2]{\newline\includegraphics[width=#1\linewidth]{#2}}
% Image côte à côte avec le texte : #1=texte, #2=frac image, #3=path, #4=url source
\newcommand{\nfimgblock}[4]{%
  \begin{minipage}[t]{0.45\linewidth}\small#1\end{minipage}%
  \hfill%
  \begin{minipage}[t]{#2\linewidth}\includegraphics[width=\linewidth]{#3}\end{minipage}%
  \source{#4}}
% Filet ornemental (séparateur de scène intra-chapitre)
\newcommand{\ornamentalrule}{%
  \vspace{\onelineskip}%
  {\centering\rule{0.4\textwidth}{0.4pt}\par}%
  \vspace{\onelineskip}%
}
% Font for scene header in notes block
\makeatletter
\ifdefined\ENDOFBOOKNOTES
  \newcommand{\scene@noteheader}[1]{{\normalfont\large #1}}
\else
  \newcommand{\scene@noteheader}[1]{{\normalfont\scshape #1}}
\fi
% \scene : public — flush inline notes (if any) then open next scene within an act.
%          Optional argument accepted but ignored (historical).
\newcommand{\scene}[1][]{%
  \ifdefined\AVECSCENES  % no-op when actes only
  \ifdefined\AVECNOTES
    \ifdefined\ENDOFBOOKNOTES\else
      \WFclear% flush pending wrapfigs before page breaks
      \cleardoublepage
      % Reset after cleardoublepage: output routine may re-corrupt these via afterpage/floats
      \global\linewidth\textwidth
      \global\@totalleftmargin\z@
      {\centering\scenotitlefont Notes\par}%
      \vspace{\onelineskip}%
      \begingroup\parindent=0pt
      \theendnotes
      \endgroup
    \fi
  \fi
  \scene@init%
  \fi
}
% \scene@notes : reset endnote counter and add scene header to notes stream.
\newcommand{\scene@notes}[1]{%
  \ifdefined\AVECNOTES
    \setcounter{endnote}{0}%
    \ifdefined\ENDOFBOOKNOTES
      \addtoendnotes{\protect\par\protect\medskip\protect\noindent\protect\scene@noteheader{#1}\protect\par\protect\smallskip}%
    \fi
  \fi
}
\makeatother

% ─── PAGE GEOMETRY (155×235 mm – Medium 1) ───────────────────────────────────
\setstocksize{235mm}{155mm}
\settrimmedsize{235mm}{155mm}{*}
\settrims{0mm}{0mm}
\setlrmarginsandblock{23mm}{18mm}{*} % inner/outer
\setulmarginsandblock{25mm}{28mm}{*} % top/bottom
\setheadfoot{11mm}{13mm}
\setheaderspaces{*}{7mm}{*}
\checkandfixthelayout

% ─── TYPOGRAPHY ───────────────────────────────────────────────────────────────
\raggedbottom
\renewcommand{\baselinestretch}{1.2}
\setlength{\parindent}{4mm}
\setlength{\parskip}{0pt}
\setlength{\afterchapskip}{0.4\onelineskip}

% ─── DIALOGUE ─────────────────────────────────────────────────────────────────
\newenvironment{dialogue}
  {\par\vspace{0.5\baselineskip}\setlength{\parindent}{0pt}}
  {\par\vspace{0.5\baselineskip}}


% ─── TABLE DES MATIÈRES ───────────────────────────────────────────────────────
\setcounter{tocdepth}{1}
\renewcommand{\chapternumberline}[1]{}           % suppress arabic number, keep Roman title only
\renewcommand{\cftchapterdotsep}{\cftdotsep}     % dot leaders
\renewcommand{\cftdot}{{\color{black!35}.}}      % lighter dots throughout TOC
\renewcommand{\cftchapterfont}{\normalfont\scshape}     % chapter entry: small caps, not bold
\renewcommand{\cftchapterpagefont}{\normalfont\scshape} % page number: small caps, not bold
% Section entries = scenes (indented, normal font)
\setlength{\cftsectionindent}{2em}
\setlength{\cftsectionnumwidth}{0em}
\renewcommand{\cftsectionfont}{\normalfont}
\renewcommand{\cftsectionpagefont}{\normalfont}
\renewcommand{\cftsectiondotsep}{\cftdotsep}
% TOC title — gallimard chapter style suppresses \printchaptertitle, so we override here
\renewcommand{\printtoctitle}[1]{{\centering\normalfont\scshape\large #1\par}}
\renewcommand{\aftertoctitle}{\par\vspace{\onelineskip}}
% Table des extraits musicaux
\newlistof{listofscores}{mus}{Table des extraits musicaux}
\newlistentry{score}{mus}{0}
\setlength{\cftscoreindent}{0pt}
\newcommand{\mylistofscores}{%
  \cleardoublepage
  \phantomsection\addcontentsline{toc}{chapter}{Table des extraits musicaux}%
  \thispagestyle{empty}%
  \vspace*{4\onelineskip}%
  {\centering\normalfont\scshape\large Table des extraits musicaux\par}%
  \vspace{2\onelineskip}%
  \begingroup
  \small\renewcommand{\baselinestretch}{1.0}\selectfont
  \listofscores*%
  \endgroup}

% \mytableofcontents : cleardoublepage + suppress self-reference entry
\newcommand{\mytableofcontents}{%
  \cleardoublepage
  \begingroup
  \setlength{\beforechapskip}{0pt}%
  \small\renewcommand{\baselinestretch}{1.0}\selectfont
  \addtocontents{toc}{\protect\setcounter{tocdepth}{-1}}%
  \tableofcontents
  \addtocontents{toc}{\protect\setcounter{tocdepth}{1}}%
  \endgroup
}

% ─── ACTE / SCENE COUNTERS ────────────────────────────────────────────────────
\newcounter{actenum}   % incremented by \acte@init
\newcounter{scenenum}  % global, incremented by \scene@init
% shared font for scene-opener number and notes title
\newcommand{\scenotitlefont}{\normalfont\large}

% ─── CHAPTER STYLE ────────────────────────────────────────────────────────────
\makeatletter
\makechapterstyle{gallimard}{%
  \renewcommand{\chapnamefont}{\normalfont\scshape\centering}
  \renewcommand{\chapnumfont}{\scenotitlefont\centering}
  \renewcommand{\chaptitlefont}{\normalfont\scshape\centering}
  \renewcommand{\printchaptername}{}
  \renewcommand{\chapternamenum}{}
  \renewcommand{\printchapternum}{%
    \ifdefined\AVECSCENES
      \chapnumfont\arabic{scenenum}%
      \markboth{\Roman{actenum}}{\arabic{scenenum}}%
    \else
      \markboth{\Roman{actenum}}{}%
    \fi}
  \renewcommand{\afterchapternum}{%
    \ifdefined\AVECSCENES\par\vskip 0.4\onelineskip\fi}
  \renewcommand{\printchaptertitle}[1]{}
}
\makeatother
\chapterstyle{gallimard}

% ─── SCENE@INIT / ACTE@INIT / ACTE@TERM ──────────────────────────────────────
\makeatletter
% \scene@openLaTeXchapter : open a LaTeX chapter hidden from TOC, then init scene notes.
\newcommand{\scene@openLaTeXchapter}[1]{%
  \addtocontents{toc}{\protect\setcounter{tocdepth}{-1}}%
  \chapter{}%
  \addtocontents{toc}{\protect\setcounter{tocdepth}{1}}%
  \scene@notes{#1}%
}
% \scene@init : increment scene counter, compute notes label, open LaTeX chapter.
\newcommand{\scene@init}{%
  \ifdefined\AVECSCENES
    \stepcounter{scenenum}%
    \ifdefined\ENDOFBOOKNOTES
      \def\scene@noteslabel{Acte~\Roman{actenum}~·~Scène~\arabic{scenenum}}%
    \else
      \def\scene@noteslabel{Acte~\Roman{actenum}~·~Scène~\arabic{scenenum}}%
    \fi
  \else
    \ifdefined\ENDOFBOOKNOTES
      \def\scene@noteslabel{Acte~\Roman{actenum}}%
    \else
      \def\scene@noteslabel{\Roman{actenum}}%
    \fi
  \fi
  \scene@openLaTeXchapter{\scene@noteslabel}%
  \ifdefined\AVECSCENES
    \addcontentsline{toc}{section}{Scène~\arabic{scenenum}}%
  \fi
}
% \acte@term : flush notes from last scene of this act (no-op when \ENDOFBOOKNOTES or \AVECSCENES absent).
\newcommand{\acte@term}{%
  \ifdefined\AVECNOTES
    \ifdefined\ENDOFBOOKNOTES\else
      \ifdefined\AVECSCENES
        \WFclear% flush pending wrapfigs before page breaks
        \cleardoublepage
        % Reset after cleardoublepage: output routine may re-corrupt these via afterpage/floats
        \global\linewidth\textwidth
        \global\@totalleftmargin\z@
        {\centering\scenotitlefont Notes\par}%
        \vspace{\onelineskip}%
        \begingroup\parindent=0pt
        \theendnotes
        \endgroup
      \fi
    \fi
  \fi
}
% \acte@init : print act heading page and open first scene.
\newcommand{\acte@init}{%
  \stepcounter{actenum}%
  \cleartorecto%
  \thispagestyle{empty}%
  \vspace*{\fill}%
  {\centering\normalfont\scshape\large Acte~\Roman{actenum}\par}%
  \vspace*{\fill}%
  \phantomsection\addcontentsline{toc}{chapter}{Acte~\Roman{actenum}}%
  \clearpage%
  \scene@init%
}
% \acte : public — full act: heading page + content + flush notes.
\newcommand{\acte}[1]{\acte@init#1\acte@term}
\makeatother

% ─── RUNNING HEADERS : verso = acte centré pc, recto = scène alignée à droite ──
\makepagestyle{gallimard}
\makeevenhead{gallimard}{}{\footnotesize\scshape\leftmark}{}
\makeoddhead{gallimard}{}{}{\footnotesize\rightmark}
\makeevenfoot{gallimard}{}{\thepage}{}
\makeoddfoot{gallimard}{}{\thepage}{}
\makepsmarks{gallimard}{%
  \nouppercaseheads
}
\pagestyle{gallimard}
\aliaspagestyle{chapter}{gallimard}

% ─── EDITORIAL NOTE (Avertissement, Postfaces) ────────────────────────────────
% 10.95pt (~small), interligne ×1.1, sans alinéa, titre en petites capitales.
\newenvironment{editorialnote}[1]{%
  \begingroup
  \leftskip=0pt\relax\rightskip=0pt\relax
  \renewcommand{\baselinestretch}{1.1}%
  \small
  \setlength{\parindent}{0pt}%
  \setlength{\leftmargini}{0pt}%
  \setlength{\parskip}{\onelineskip}%
  \ifx&#1&\else
    {\centering\normalfont\scshape\large #1\par}%
    \vspace{\onelineskip}%
  \fi
  \noindent\ignorespaces
}{%
  \par\endgroup
}

%   \begin{document}
%   ...
% ─────────────────────────────────────────────────────────────────────────────

% ─── PDF VERSION ──────────────────────────────────────────────────────────────
\ifdefined\pdfvariable\pdfvariable minorversion 7\fi

% ─── PACKAGES ─────────────────────────────────────────────────────────────────
\usepackage{fontspec}
% ────── EB Garamond
%%% \setmainfont{EB Garamond}
% ────── Cormorant Garamond
%%% \setmainfont[
%%%   UprightFont    = *-Regular,
%%%   ItalicFont     = *-Italic,
%%%   BoldFont       = *-SemiBold,
%%%   BoldItalicFont = *-SemiBoldItalic
%%% ]{Cormorant Garamond}
% ────── EB Garamond (main) + Adobe Garamond Pro (bold only)
\setmainfont{EB Garamond}[
  BoldFont       = AGaramondPro-Semibold,
  BoldItalicFont = AGaramondPro-SemiboldItalic,
]

\newfontfamily{\arabicfont}{Geeza Pro}
\newcommand{\arab}[1]{{\arabicfont\textdir TRT #1}}
\newfontfamily{\cjkfont}{Hiragino Mincho ProN}
\newcommand{\cjk}[1]{{\cjkfont #1}}
\newcommand{\shouting}[1]{\textbf{\textsc{#1}}}
\usepackage[russian,french]{babel}
\usepackage[program=/usr/local/bin/lilypond,includepaths={music}]{lyluatex}
\usepackage{microtype}
\usepackage{fmtcount}
\usepackage{catchfile}
\usepackage{graphicx}
\usepackage{xcolor}
\usepackage{amsmath}
\usepackage{xparse}
\usepackage{tikz}
\usepackage{afterpage}
\usepackage{needspace}
\usepackage{pdfpages}
\usepackage[absolute,overlay]{textpos}
\usepackage{tabularx}
\usepackage{booktabs}
\usepackage{calc}
\usepackage{wrapfig}
\usepackage{tcolorbox}
\usepackage{ragged2e}
% Fix: \WF@finale restores \linewidth and \@totalleftmargin only when
% \parshape==\WF@fudgeparshape; the "forced to float" collision path bypasses
% that guard, leaving both corrupted globally.  Patch \WF@finale to force-reset
% both after every call, regardless of which path was taken.
\makeatletter
\let\WF@@orig@finale\WF@finale
\def\WF@finale{\WF@@orig@finale\global\linewidth\textwidth\global\@totalleftmargin\z@}
\makeatother
% Safety net: clear any residual wrapfig state at chapter boundaries.
\pretocmd{\chapter}{\WFclear}{}{}
\usepackage{hyperref}
\usepackage{bookmark}
\hypersetup{pdfpagelayout=TwoPageRight, hidelinks}
\usetikzlibrary{calc,positioning}


% ─── VERSION / GIT ────────────────────────────────────────────────────────────
\newcommand{\docversion}{0.5.0}
\newcommand{\docversionordinal}{Deuxième}

\IfFileExists{tete_de_veau_ravigote.gitinfo}%
  {\CatchFileDef{\gitcommit}{tete_de_veau_ravigote.gitinfo}{\endlinechar=-1}}%
  {\def\gitcommit{}}

% ─── GRAYSCALE IMAGE FILTER ───────────────────────────────────────────────────
\graphicspath{{iconography/}}
\newcommand{\figcaptionname}{\footnotesize}
\newcommand{\figcaptiondetail}{\scriptsize}
\newcommand{\figpagetitlename}{\small\scshape}
\newcommand{\figpagetitledetail}{\footnotesize}
\newcommand{\bwimage}[2][]{%
  \begin{tikzpicture}
    \node[inner sep=0, outer sep=0pt, fill=white] (img) {\includegraphics[#1]{#2}};
    %%% begin{scope}[blend mode=saturation]
    %%%  \fill[black] (img.south west) rectangle (img.north east);
    %%% end{scope}
  \end{tikzpicture}%
}

% ─── ESPACEMENT FLOTTANTS ─────────────────────────────────────────────────────
\setlength{\textfloatsep}{8pt}   % espace entre flottant top/bottom et le texte (défaut ~20pt)

% ─── ICONOGRAPHIE ─────────────────────────────────────────────────────────────
% All commands below are no-ops unless \AVECILLUSTRATIONS is defined.
% Image paths are relative to the graphicspath (iconography/).
\ifdefined\AVECILLUSTRATIONS
  \newlength{\iconofillheight}% scratch register for [p] page-fill placement
  % \iconographieimg[scale]{image}{caption}
  % Single image centred on its own dedicated page (via \afterpage).
  % [scale] height as fraction of \textheight (default 0.75); {caption} may be empty.
  \newcommand{\iconographieimg}[3][0.75]{%
    \afterpage{%
      \thispagestyle{gallimard}%
      \vspace*{\fill}%
      \begin{center}%
        \bwimage[height=#1\textheight,width=\textwidth,keepaspectratio]{#2}\\[0.6em]%
        {\figcaptionname\itshape #3}%
      \end{center}%
      \vspace*{\fill}%
      \newpage
    }%
  }
  % \iconographietex{figures/foo_fig.tex}
  % Inserts one hand-crafted figure page from an external .tex file (via \afterpage).
  % The file is responsible for its own layout — see figures/*_fig.tex for examples.
  \newcommand{\iconographietex}[1]{%
    \afterpage{%
      \thispagestyle{gallimard}%
      \input{#1}%
      \clearpage
    }%
  }
  % \iconographietexpair{figures/foo_fig.tex}{figures/bar_fig.tex}
  % Like \iconographietex but inserts two consecutive figure pages.
  \newcommand{\iconographietexpair}[2]{%
    \afterpage{%
      \thispagestyle{gallimard}%
      \input{#1}%
      \input{#2}%
      \clearpage
    }%
  }
  % \iconographiewrapfig{pos}{width}{\bwimage[width=\linewidth]{path}}{caption}
  % Image wrapped by surrounding text (wrapfigure).
  % {pos} = l or r;  {width} = e.g. 0.35\textwidth;  {caption} may be empty {}.
  \newcommand{\iconographiewrapfig}[4]{%
    \setlength{\intextsep}{0pt}%
    \setlength{\columnsep}{8pt}%
    \begin{wrapfigure}{#1}{#2}%
      \centering
      #3%
      \ifx&#4&\else\\[0.2em]{\figcaptionname\itshape #4}\\[0.4\baselineskip]\fi%
    \end{wrapfigure}%
  }
  % \iconographieinlineblock[placement][lines]{\bwimage[width=W]{path}}{caption}
  % Centred image block in the text flow. {caption} may be empty {}.
  % [placement] empty or omitted → inline: caller must put \par before and \noindent after.
  % [placement] = t or b → LaTeX float (top/bottom of page); no \par/\noindent needed.
  % [placement] = p → fill remaining page height inline (image scaled to fit, no float).
  % [lines] number of caption lines (default 1); used by [p] to reserve the right height.
  % Multiple inline calls in the same paragraph are supported.
  \NewDocumentCommand{\iconographieinlineblock}{O{} O{1} m m}{%
    % #1=placement  #2=caption lines (for [p])  #3=image  #4=caption
    \ifx&#1&%
      % no placement → inline
      \vspace{0.5\onelineskip}%
      \centerline{#3}%
      \ifx&#4&\else\par\centerline{\figcaptionname\itshape#4}\vspace{0.5\baselineskip}\fi%
      \vspace{-0.5\onelineskip}%
    \else
      \ifthenelse{\equal{#1}{p}}{%
        % p → scale image to fill remaining page height
        \par
        \begingroup
          \setlength{\iconofillheight}{\dimexpr\pagegoal-\pagetotal\relax}%
          \addtolength{\iconofillheight}{-2\baselineskip}%
          \ifx&#4&\else\addtolength{\iconofillheight}{-#2\baselineskip-0.5\baselineskip}\fi%
          \ifdim\iconofillheight>0pt
            \vspace{2\baselineskip}%
            \centerline{\resizebox*{!}{\iconofillheight}{#3}}%
            \ifx&#4&\else\par{\centering\figcaptionname\itshape#4\par}\fi%
          \fi
        \endgroup
      }{%
        % placement given → float
        \begin{figure}[#1]
          \centering#3%
          \ifx&#4&\else\par{\figcaptionname\itshape#4}\fi%
        \end{figure}%
      }%
    \fi
  }
  % \iconographiedouble[placement][height][margin]{\bwimage{img1}}{\bwimage{img2}}[caption]
  % Two images side by side at equal height.
  % [placement] = b (default, push to bottom) or t (push to top)
  % [height] common image height (default 5cm); [margin] symmetric indent (default 0pt)
  % [caption] shared caption, omit for none.
  \NewDocumentCommand{\iconographiedouble}{O{b} O{5cm} O{0pt} m m O{}}{%
    \ifthenelse{\equal{#1}{b}}{\vspace*{\fill}}{\vspace*{0pt}}%
    \vspace{0.5\onelineskip}%
    {\setlength{\leftskip}{#3}\setlength{\rightskip}{#3}%
    \noindent\bwimage[height=#2]{#4}\hfill\bwimage[height=#2]{#5}\par}%
    \ifx&#6&%
      \vspace{0.5\onelineskip}%
    \else%
      \vspace{0.2\onelineskip}%
      {\centering\figcaptionname\itshape#6\par}%
      \vspace{0.3\onelineskip}%
    \fi%
    \ifthenelse{\equal{#1}{t}}{\vspace*{\fill}}{}%
  }
  % \iconographietriple[placement][height][margin]{\bwimage{img1}}{\bwimage{img2}}{\bwimage{img3}}[caption]
  % Three images side by side at equal height. Same options as \iconographiedouble.
  \NewDocumentCommand{\iconographietriple}{O{b} O{5cm} O{0pt} m m m O{}}{%
    \ifthenelse{\equal{#1}{b}}{\vspace*{\fill}}{\vspace*{0pt}}%
    \vspace{0.5\onelineskip}%
    {\setlength{\leftskip}{#3}\setlength{\rightskip}{#3}%
    \noindent\bwimage[height=#2]{#4}\hfill\bwimage[height=#2]{#5}\hfill\bwimage[height=#2]{#6}\par}%
    \ifx&#7&%
      \vspace{0.5\onelineskip}%
    \else%
      \vspace{0.2\onelineskip}%
      {\centering\figcaptionname\itshape#7\par}%
      \vspace{0.3\onelineskip}%
    \fi%
    \ifthenelse{\equal{#1}{t}}{\vspace*{\fill}}{}%
  }
\else
  \newcommand{\iconographieimg}[3][0.75]{}
  \newcommand{\iconographietex}[1]{}
  \newcommand{\iconographietexpair}[2]{}
  \newcommand{\iconographiewrapfig}[4]{}
  \NewDocumentCommand{\iconographieinlineblock}{O{} O{1} m m}{}
  \NewDocumentCommand{\iconographiedouble}{O{b} O{5cm} O{0pt} m m O{}}{}
  \NewDocumentCommand{\iconographietriple}{O{b} O{5cm} O{0pt} m m m O{}}{}
\fi

% ─── EXTRAITS MUSICAUX ────────────────────────────────────────────────────────
% Compiler avec : lualatex "\def\AVECMUSIQUE{}\input{...}"
\ifdefined\AVECMUSIQUE
  % Extrait musical inline : [#1]=espacement (défaut 0.66), #2=fichier dans music/, #3=légende
  \newcommand{\score}[3][0.66]{%
    \addcontentsline{mus}{score}{#3}%
    \noindent\begin{minipage}{\textwidth}%
    \vspace{#1\onelineskip}%
    \lilypondfile[line-width=\textwidth]{music/#2}%
    \begin{center}%
      \vspace{-#1\onelineskip}%
      {\centering\figcaptionname\itshape#3}%
      \vspace{#1\onelineskip}%
    \end{center}%
    \end{minipage}}
\else
  \newcommand{\score}[3][1]{}
\fi

% ─── NOTES FINALES ────────────────────────────────────────────────────────────
% Compiler avec : lualatex "\def\AVECNOTES{}\input{...}"
% \ENDOFBOOKNOTES : notes regroupées en fin de livre (avec en-tête de scène) ;
%                   absent = notes à la fin de chaque scène (sans en-tête).
\usepackage{endnotes}
\ifdefined\AVECNOTES
  \newcommand{\nf}[1]{\endnote{#1}}
\else
  \newcommand{\nf}[1]{}
\fi
% Supprimer le titre automatique du package (on gère le nôtre)
\renewcommand{\notesname}{}
% Style des notes : même fonte, corps légèrement réduit
\renewcommand{\enotesize}{\small}
% Format des notes en fin de livre : sans retrait, numéro en corps normal
\makeatletter
\renewcommand\enoteformat{%
  \rightskip\z@ \leftskip\z@ \parindent=0pt
  \leavevmode\theenmark.\enspace}
\makeatother
% Fix PDF ActualText for œ (oe-ligature) copy-paste in \source URLs
\usepackage{accsupp}
\newcommand{\oechar}{\char"153\relax}
{\catcode"153=\active\gdefœ{\BeginAccSupp{method=hex,unicode,ActualText=0153}\oechar\EndAccSupp{}}}
% Source dans les notes : saut de ligne avant le lien Wikipedia
\newcommand{\source}[1]{\newline\textit{Source~:}~{\upshape\addfontfeatures{RawFeature=-tlig;-calt}\def\_{\textunderscore\allowbreak}\def\%{\char"25\relax\allowbreak}\catcode"153=\active\scantokens{#1\empty}}}
% Image seule dans une note (centrée sous le texte)
\newcommand{\nfimg}[2]{\newline\includegraphics[width=#1\linewidth]{#2}}
% Image côte à côte avec le texte : #1=texte, #2=frac image, #3=path, #4=url source
\newcommand{\nfimgblock}[4]{%
  \begin{minipage}[t]{0.45\linewidth}\small#1\end{minipage}%
  \hfill%
  \begin{minipage}[t]{#2\linewidth}\includegraphics[width=\linewidth]{#3}\end{minipage}%
  \source{#4}}
% Filet ornemental (séparateur de scène intra-chapitre)
\newcommand{\ornamentalrule}{%
  \vspace{\onelineskip}%
  {\centering\rule{0.4\textwidth}{0.4pt}\par}%
  \vspace{\onelineskip}%
}
% Font for scene header in notes block
\makeatletter
\ifdefined\ENDOFBOOKNOTES
  \newcommand{\scene@noteheader}[1]{{\normalfont\large #1}}
\else
  \newcommand{\scene@noteheader}[1]{{\normalfont\scshape #1}}
\fi
% \scene : public — flush inline notes (if any) then open next scene within an act.
%          Optional argument accepted but ignored (historical).
\newcommand{\scene}[1][]{%
  \ifdefined\AVECSCENES  % no-op when actes only
  \ifdefined\AVECNOTES
    \ifdefined\ENDOFBOOKNOTES\else
      \WFclear% flush pending wrapfigs before page breaks
      \cleardoublepage
      % Reset after cleardoublepage: output routine may re-corrupt these via afterpage/floats
      \global\linewidth\textwidth
      \global\@totalleftmargin\z@
      {\centering\scenotitlefont Notes\par}%
      \vspace{\onelineskip}%
      \begingroup\parindent=0pt
      \theendnotes
      \endgroup
    \fi
  \fi
  \scene@init%
  \fi
}
% \scene@notes : reset endnote counter and add scene header to notes stream.
\newcommand{\scene@notes}[1]{%
  \ifdefined\AVECNOTES
    \setcounter{endnote}{0}%
    \ifdefined\ENDOFBOOKNOTES
      \addtoendnotes{\protect\par\protect\medskip\protect\noindent\protect\scene@noteheader{#1}\protect\par\protect\smallskip}%
    \fi
  \fi
}
\makeatother

% ─── PAGE GEOMETRY (155×235 mm – Medium 1) ───────────────────────────────────
\setstocksize{235mm}{155mm}
\settrimmedsize{235mm}{155mm}{*}
\settrims{0mm}{0mm}
\setlrmarginsandblock{23mm}{18mm}{*} % inner/outer
\setulmarginsandblock{25mm}{28mm}{*} % top/bottom
\setheadfoot{11mm}{13mm}
\setheaderspaces{*}{7mm}{*}
\checkandfixthelayout

% ─── TYPOGRAPHY ───────────────────────────────────────────────────────────────
\raggedbottom
\renewcommand{\baselinestretch}{1.2}
\setlength{\parindent}{4mm}
\setlength{\parskip}{0pt}
\setlength{\afterchapskip}{0.4\onelineskip}

% ─── DIALOGUE ─────────────────────────────────────────────────────────────────
\newenvironment{dialogue}
  {\par\vspace{0.5\baselineskip}\setlength{\parindent}{0pt}}
  {\par\vspace{0.5\baselineskip}}


% ─── TABLE DES MATIÈRES ───────────────────────────────────────────────────────
\setcounter{tocdepth}{1}
\renewcommand{\chapternumberline}[1]{}           % suppress arabic number, keep Roman title only
\renewcommand{\cftchapterdotsep}{\cftdotsep}     % dot leaders
\renewcommand{\cftdot}{{\color{black!35}.}}      % lighter dots throughout TOC
\renewcommand{\cftchapterfont}{\normalfont\scshape}     % chapter entry: small caps, not bold
\renewcommand{\cftchapterpagefont}{\normalfont\scshape} % page number: small caps, not bold
% Section entries = scenes (indented, normal font)
\setlength{\cftsectionindent}{2em}
\setlength{\cftsectionnumwidth}{0em}
\renewcommand{\cftsectionfont}{\normalfont}
\renewcommand{\cftsectionpagefont}{\normalfont}
\renewcommand{\cftsectiondotsep}{\cftdotsep}
% TOC title — gallimard chapter style suppresses \printchaptertitle, so we override here
\renewcommand{\printtoctitle}[1]{{\centering\normalfont\scshape\large #1\par}}
\renewcommand{\aftertoctitle}{\par\vspace{\onelineskip}}
% Table des extraits musicaux
\newlistof{listofscores}{mus}{Table des extraits musicaux}
\newlistentry{score}{mus}{0}
\setlength{\cftscoreindent}{0pt}
\newcommand{\mylistofscores}{%
  \cleardoublepage
  \phantomsection\addcontentsline{toc}{chapter}{Table des extraits musicaux}%
  \thispagestyle{empty}%
  \vspace*{4\onelineskip}%
  {\centering\normalfont\scshape\large Table des extraits musicaux\par}%
  \vspace{2\onelineskip}%
  \begingroup
  \small\renewcommand{\baselinestretch}{1.0}\selectfont
  \listofscores*%
  \endgroup}

% \mytableofcontents : cleardoublepage + suppress self-reference entry
\newcommand{\mytableofcontents}{%
  \cleardoublepage
  \begingroup
  \setlength{\beforechapskip}{0pt}%
  \small\renewcommand{\baselinestretch}{1.0}\selectfont
  \addtocontents{toc}{\protect\setcounter{tocdepth}{-1}}%
  \tableofcontents
  \addtocontents{toc}{\protect\setcounter{tocdepth}{1}}%
  \endgroup
}

% ─── ACTE / SCENE COUNTERS ────────────────────────────────────────────────────
\newcounter{actenum}   % incremented by \acte@init
\newcounter{scenenum}  % global, incremented by \scene@init
% shared font for scene-opener number and notes title
\newcommand{\scenotitlefont}{\normalfont\large}

% ─── CHAPTER STYLE ────────────────────────────────────────────────────────────
\makeatletter
\makechapterstyle{gallimard}{%
  \renewcommand{\chapnamefont}{\normalfont\scshape\centering}
  \renewcommand{\chapnumfont}{\scenotitlefont\centering}
  \renewcommand{\chaptitlefont}{\normalfont\scshape\centering}
  \renewcommand{\printchaptername}{}
  \renewcommand{\chapternamenum}{}
  \renewcommand{\printchapternum}{%
    \ifdefined\AVECSCENES
      \chapnumfont\arabic{scenenum}%
      \markboth{\Roman{actenum}}{\arabic{scenenum}}%
    \else
      \markboth{\Roman{actenum}}{}%
    \fi}
  \renewcommand{\afterchapternum}{%
    \ifdefined\AVECSCENES\par\vskip 0.4\onelineskip\fi}
  \renewcommand{\printchaptertitle}[1]{}
}
\makeatother
\chapterstyle{gallimard}

% ─── SCENE@INIT / ACTE@INIT / ACTE@TERM ──────────────────────────────────────
\makeatletter
% \scene@openLaTeXchapter : open a LaTeX chapter hidden from TOC, then init scene notes.
\newcommand{\scene@openLaTeXchapter}[1]{%
  \addtocontents{toc}{\protect\setcounter{tocdepth}{-1}}%
  \chapter{}%
  \addtocontents{toc}{\protect\setcounter{tocdepth}{1}}%
  \scene@notes{#1}%
}
% \scene@init : increment scene counter, compute notes label, open LaTeX chapter.
\newcommand{\scene@init}{%
  \ifdefined\AVECSCENES
    \stepcounter{scenenum}%
    \ifdefined\ENDOFBOOKNOTES
      \def\scene@noteslabel{Acte~\Roman{actenum}~·~Scène~\arabic{scenenum}}%
    \else
      \def\scene@noteslabel{Acte~\Roman{actenum}~·~Scène~\arabic{scenenum}}%
    \fi
  \else
    \ifdefined\ENDOFBOOKNOTES
      \def\scene@noteslabel{Acte~\Roman{actenum}}%
    \else
      \def\scene@noteslabel{\Roman{actenum}}%
    \fi
  \fi
  \scene@openLaTeXchapter{\scene@noteslabel}%
  \ifdefined\AVECSCENES
    \addcontentsline{toc}{section}{Scène~\arabic{scenenum}}%
  \fi
}
% \acte@term : flush notes from last scene of this act (no-op when \ENDOFBOOKNOTES or \AVECSCENES absent).
\newcommand{\acte@term}{%
  \ifdefined\AVECNOTES
    \ifdefined\ENDOFBOOKNOTES\else
      \ifdefined\AVECSCENES
        \WFclear% flush pending wrapfigs before page breaks
        \cleardoublepage
        % Reset after cleardoublepage: output routine may re-corrupt these via afterpage/floats
        \global\linewidth\textwidth
        \global\@totalleftmargin\z@
        {\centering\scenotitlefont Notes\par}%
        \vspace{\onelineskip}%
        \begingroup\parindent=0pt
        \theendnotes
        \endgroup
      \fi
    \fi
  \fi
}
% \acte@init : print act heading page and open first scene.
\newcommand{\acte@init}{%
  \stepcounter{actenum}%
  \cleartorecto%
  \thispagestyle{empty}%
  \vspace*{\fill}%
  {\centering\normalfont\scshape\large Acte~\Roman{actenum}\par}%
  \vspace*{\fill}%
  \phantomsection\addcontentsline{toc}{chapter}{Acte~\Roman{actenum}}%
  \clearpage%
  \scene@init%
}
% \acte : public — full act: heading page + content + flush notes.
\newcommand{\acte}[1]{\acte@init#1\acte@term}
\makeatother

% ─── RUNNING HEADERS : verso = acte centré pc, recto = scène alignée à droite ──
\makepagestyle{gallimard}
\makeevenhead{gallimard}{}{\footnotesize\scshape\leftmark}{}
\makeoddhead{gallimard}{}{}{\footnotesize\rightmark}
\makeevenfoot{gallimard}{}{\thepage}{}
\makeoddfoot{gallimard}{}{\thepage}{}
\makepsmarks{gallimard}{%
  \nouppercaseheads
}
\pagestyle{gallimard}
\aliaspagestyle{chapter}{gallimard}

% ─── EDITORIAL NOTE (Avertissement, Postfaces) ────────────────────────────────
% 10.95pt (~small), interligne ×1.1, sans alinéa, titre en petites capitales.
\newenvironment{editorialnote}[1]{%
  \begingroup
  \leftskip=0pt\relax\rightskip=0pt\relax
  \renewcommand{\baselinestretch}{1.1}%
  \small
  \setlength{\parindent}{0pt}%
  \setlength{\leftmargini}{0pt}%
  \setlength{\parskip}{\onelineskip}%
  \ifx&#1&\else
    {\centering\normalfont\scshape\large #1\par}%
    \vspace{\onelineskip}%
  \fi
  \noindent\ignorespaces
}{%
  \par\endgroup
}
